% Part: proof-theory
% Chapter: natural-deduction
% Section: sequents

\documentclass[../../../include/open-logic-section]{subfiles}

\begin{document}

\olfileid{pt}{nat}{seq}

\olsection{في الاستنباط الطبيعي بالمتتاليات: \Log{N2c} و\Log{N2i}}

يدل~!!^a{proof}~$\delta$ في~\Log{N1c} أو~\Log{N1i} على أن~!!{formula}
التي في نهايته، وهي~$!A$، تلزم من فروضه المفتوحة. فإذا كانت~$\Gamma$
مجموعة~!!{formula}s~$!B$ التي يقع منها~$!B^x$ فرضًا مفتوحًا في~$\delta$،
رمزنا إلى ذلك بـ~$\delta \Proves \Gamma \Sequent !A$. ومن هنا يستبين
إمكان بناء استنباط طبيعي يعمل في المتتاليات. فكل متتالية في هذا النظام
تبين الفروض المفتوحة التي تعتمد عليها~!!{formula} الواقعة عن اليمين،
أي فروض~!!{proof} الجزئي المنتهي بتلك المتتالية. وهذا هو الاستنباط
الطبيعي بالمتتاليات، أو الاستنباط الطبيعي «على طريقة المتتاليات»،
ونسمي النظامين الحاصلين~\Log{N2c} و\Log{N2i}.

ولكي تقرب المطابقة بين~\Log{N1} و\Log{N2}، نجعل المجموعات~$\Gamma$
الواقعة عن اليسار مؤلفة من~!!{formula}s موسومة~$x:!B$. وقواعد~\Log{N1}
لا تعمل إلا في~!!{formula}s المقدمات، أي~!!{formula}s التي تنتهي بها
!!{proof}s الجزئية المباشرة؛ فنجعل قواعد~\Log{N2} لا تعمل إلا في تالي
المتتالية. وأما الصيغ الموسومة عن يسار المتتاليات فنسميها
\emph{السياق}.

وبإزاء الفروض تكون المتتاليات العليا في براهين~\Log{N2} مسلّمات من الصورة
\[x:!A \Sequent !A.\]
ثم تطابق القواعد قواعد~\Log{N1} على التمام، وهي مثبتة في~\olref{tab:N2}.

\subfile{rules-N2}

وكما أجزنا في~\Log{N1c} و\Log{N1i} لقواعد~$\Intro{\lif}$ و
$\Intro{\lnot}$ و$\Elim{\lor}$ و$\Elim{\lexists}$ أن تبرئ فروضًا أو
لا تبرئها، نجيز التطبيقات الصحيحة للقواعد المقابلة في~\Log{N2c}
و\Log{N2i}، ولو لم تقع~!!{formula}s السياق المقابلة~$!A$ و$!B$
و$!A(a)$ في سياق المقدمة المعنية.

\begin{prop}
إذا كان~$\delta$~!!a{proof} في~\Log{N1c} أو~\Log{N1i} لـ$!A$، وُجد
!!a{proof}~$\delta'$ لـ$\Gamma \Sequent !A$ في~\Log{N2c}، أو في
\Log{N2i}، حيث~$\Gamma$ مجموعة كل~$x:!B$ التي يكون فيها~$!B^x$ فرضًا
مفتوحًا في~$\delta$.
\end{prop}

\begin{proof}
  نسلك الاستقراء على~$n=\pheight{\delta}$. فإذا كان~$n=0$، لم يكن
  $\delta$ إلا الفرض المفتوح~$!A^x$. وعندئذ~$\Gamma=\{x:!A\}$، وتكون
  $\Gamma \Sequent !A$ مسلّمة في~\Log{N2c} أو~\Log{N2i}.

  وإن كان~$n>0$ قسمنا بحسب آخر الاستدلال في~$\delta$؛ ونبين حالة واحدة،
  ونكل سائر الحالات إلى التمرين.

  لنفرض أن آخر الاستدلال~$\Intro{\lif}$. عندئذ ينتهي~$\delta$ بـ
  \[
  \AxiomC{$\Discharge{!B}{x}$}
  \RightLabel{$\delta_1$}
  \DeduceC{$!C$}
  \DischargeRule{\Intro{\lif}}{x}
  \UnaryInfC{$!B \lif !C$}
  \DisplayProof
  \]
  وبفرض الاستقراء يوجد~!!a{proof}~$\delta_1'$ في~\Log{N2c}، أو في
  \Log{N2i}، لـ$\Gamma' \Sequent !C$، حيث~$\Gamma'$ هي الفروض الموسومة
  المفتوحة في~$\delta_1$. فإذا كان~$!B^x$ فرضًا مفتوحًا في~$\delta_1$،
  بُرئ عند استدلال~$\Intro{\lif}$، وكانت فروض~$\delta$ المفتوحة هي
  $\Gamma=\Gamma'\setminus\{x:!B\}$. وبعبارة أخرى، يكون~$\delta_1'$
  !!a{proof} لـ$x:!B, \Gamma \Sequent !C$، ويمكننا اتخاذ~$\delta'$ هو
  \[
  \AxiomC{}
  \RightLabel{$\delta_1'$}
  \Deduce$x:!B, \Gamma \fCenter !C$
  \RightLabel{\Intro{\lif}}
  \UnaryInf$\Gamma \fCenter !B \lif !C$
  \DisplayProof
  \]
  وإذا لم يكن~$!B^x$ فرضًا مفتوحًا في~$\delta_1$، لم يبرئ
  $\Intro{\lif}$ فرضًا، وكانت الفروض المفتوحة في~$\delta$ و$\delta_1$
  واحدة. ولأن وجود صيغة السياق~$x:!B$ في مقدمة~$\Intro{\lif}$
  ليس لازمًا في~\Log{N2c} و\Log{N2i}، أمكننا اتخاذ~$\delta'$ هو
  \[
  \AxiomC{}
  \RightLabel{$\delta_1'$}
  \Deduce$\Gamma \fCenter !C$
  \RightLabel{\Intro{\lif}}
  \UnaryInf$\Gamma \fCenter !B \lif !C$
  \DisplayProof
  \]

  والحالات التي ينتهي فيها~$\delta$ بقاعدة أخرى مماثلة.
\end{proof}

\begin{prop}
إذا كان~$\delta$~!!a{proof} في~\Log{N2c} أو~\Log{N2i} لـ
$\Gamma \Sequent !A$، وُجد~!!a{proof}~$\delta'$ في~\Log{N1c}، أو في
\Log{N1i}، تكون~!!{formula} نهايته~$!A$، وفروضه المفتوحة~$!B^x$ لكل
$x:!B \in \Gamma$.
\end{prop}

\begin{proof}
  ونسلك أيضًا الاستقراء على~!!{height}~$n$ للـ$\delta$. فإذا كان
  $n=0$، كان~$\delta$ المسلّمة~$x:!A \Sequent !A$. ويكون~!!{proof}
  $\delta'$ هو الفرض~$!A^x$.

  وإن كان~$n>0$ قسمنا بحسب آخر الاستدلال في~$\delta$. فلو كان ذلك
  الاستدلال~$\Intro{\lif}$ مثلًا، انتهى~$\delta$ بـ
  \[
  \bottomAlignProof
  \AxiomC{}
  \RightLabel{$\delta_1$}
  \Deduce$x:!B, \Gamma \fCenter !C$
  \RightLabel{\Intro{\lif}}
  \UnaryInf$\Gamma \fCenter !B \lif !C$
  \DisplayProof
  \quad\text{ أو }\quad
  \bottomAlignProof
  \AxiomC{}
  \RightLabel{$\delta_1$}
  \Deduce$\Gamma \fCenter !C$
  \RightLabel{\Intro{\lif}}
  \UnaryInf$\Gamma \fCenter !B \lif !C$
  \DisplayProof
  \]
  وبفرض الاستقراء يوجد~!!a{proof}~$\delta_1'$ في~\Log{N1c}، أو في
  \Log{N1i}، لـ$!C$. وفي الحالة الأولى يكون~$!B^x$ فرضًا مفتوحًا في
  $\delta_1'$، ولا يكون كذلك في الحالة الثانية. ويمكننا اتخاذ
  !!{proof}~$\delta'$ هو، على الترتيب،
  \[
  \bottomAlignProof
  \AxiomC{$\Discharge{!B}{x}$}
  \RightLabel{$\delta_1'$}
  \DeduceC{$!C$}
  \DischargeRule{\Intro{\lif}}{x}
  \UnaryInfC{$!B \lif !C$}
  \DisplayProof
  \quad\text{ أو }\quad
  \bottomAlignProof
  \AxiomC{}
  \RightLabel{$\delta_1'$}
  \DeduceC{$!C$}
  \RightLabel{\Intro{\lif}}
  \UnaryInfC{$!B \lif !C$}
  \DisplayProof
  \]
\end{proof}

\end{document}
